\makeatletter
\def\input@path{{./}{../}{../../}{../../../}{../../../../}{preamble/}{../preamble/}{../../preamble/}{../../../preamble/}}
\makeatother

%\input{preamble_G10_T1_practice_activities_c1_c3.tex}
\input{preamble.tex}
\input{graphs/highlights.tex}
% \GTenHighlightsfalse % Uncomment to disable instructional highlights.
\input{graphs/specific_highlights.tex}


\title{10th Grade Midterm Analysis of Decisions: C4 Solutions}
\author{}
\date{}

\pagestyle{empty}

\begin{document}
	
	\maketitle
	\setcounter{tocdepth}{1}
	\setcounter{secnumdepth}{2}
	\phantomsection\label{toc}
	\vspace*{-20mm}
	\tableofcontents
	
	% ============================================================
	% CRITERION
	% ============================================================
	
	\section*{Evaluated Criterion}
	
	C6: Calculates Expected Monetary Value using stated or estimated probabilities and applies the Bayes decision criterion to select the alternative with the highest expected payoff.
	
	% ============================================================
	% INTRODUCTION
	% ============================================================
	
	\section{Introduction}\label{sec:c4-introduction}
\phantomsection\label{tone:legend}
\PracticeHighlightLegend{G10-T1-C6}

	\SectionProblemsTableOfContents{
		\SectionProblemLink{sec:c4-intro-purpose}{Expected Value Criterion} \\
		\SectionProblemLink{sec:c4-intro-2x2}{General 2x2 Problem} \\
		\SectionProblemLink{sec:c4-intro-4x3}{General 4x3 Problem}
	}\vspace{8mm}
	
	\subsection{Expected Value Criterion}\label{sec:c4-intro-purpose}

	
	When \GTenToneHL{probability}{probabilities} can be assigned to the possible \GTenToneHL{state}{states} \GTenToneHL{state}{of} \GTenToneHL{state}{nature},
	a decision-maker can compare \GTenToneHL{alternative}{alternatives} using their expected monetary
	values. The \GTenToneHL{expected}{Expected} \GTenToneHL{expected}{Monetary} \GTenToneHL{expected}{Value} (\GTenToneHL{expected}{EMV}) criterion, also known as the Bayes decision criterion, combines the \GTenToneHL{payoff}{payoff} associated with each possible \GTenToneHL{state}{state} \GTenToneHL{state}{of} \GTenToneHL{state}{nature} with the \GTenToneHL{probability}{probability} that the state will occur.
	For a \GTenToneHL{payoff}{payoff} table with \GTenToneHL{alternative}{alternatives} in rows and \GTenToneHL{state}{states} \GTenToneHL{state}{of} \GTenToneHL{state}{nature} in
	columns, let
	$
	\GTenToneHL{payoff}{P_{ay}}(\GTenToneHL{alternative}{A_i},\GTenToneHL{state}{S_j})
	$
	represent the \GTenToneHL{payoff}{payoff} obtained when \GTenToneHL{alternative}{alternative} $\GTenToneHL{alternative}{A_i}$ is selected and
	\GTenToneHL{state}{state} \GTenToneHL{state}{of} \GTenToneHL{state}{nature} $\GTenToneHL{state}{S_j}$ occurs.
	%
	\begin{center}
		Let $\GTenToneHL{probability}{P}(\GTenToneHL{state}{S_j})$ represent the \GTenToneHL{probability}{probability} that \GTenToneHL{state}{state} \GTenToneHL{state}{of} \GTenToneHL{state}{nature} $\GTenToneHL{state}{S_j}$ occurs.
	\end{center}
	%
	The \GTenToneHL{probability}{probabilities} must satisfy
	$$
	\sum_{j=1}^{n} \GTenToneHL{probability}{P}(\GTenToneHL{state}{S_j}) = 1.
	$$
	%
	For each \GTenToneHL{alternative}{alternative}, calculate its \GTenToneHL{expected}{Expected} \GTenToneHL{expected}{Monetary} \GTenToneHL{expected}{Value}:
	$$
	\GTenToneHL{expected}{EMV}(\GTenToneHL{alternative}{A_i})
	=
	\sum_{j=1}^{n}
	\GTenToneHL{probability}{P}(\GTenToneHL{state}{S_j})\GTenToneHL{payoff}{P_{ay}}(\GTenToneHL{alternative}{A_i},\GTenToneHL{state}{S_j})
	$$
	
	or, explicitly,
	$$
	\GTenToneHL{expected}{EMV}(\GTenToneHL{alternative}{A_i})
	=
	\GTenToneHL{probability}{P}(\GTenToneHL{state}{S_1})\GTenToneHL{payoff}{P_{ay}}(\GTenToneHL{alternative}{A_i},\GTenToneHL{state}{S_1})
	+
	\GTenToneHL{probability}{P}(\GTenToneHL{state}{S_2})\GTenToneHL{payoff}{P_{ay}}(\GTenToneHL{alternative}{A_i},\GTenToneHL{state}{S_2})
	+
	\cdots
	+
	\GTenToneHL{probability}{P}(\GTenToneHL{state}{S_n})\GTenToneHL{payoff}{P_{ay}}(\GTenToneHL{alternative}{A_i},\GTenToneHL{state}{S_n}).
	$$
	
	The Bayes decision is the \GTenToneHL{alternative}{alternative} with the highest expected \GTenToneHL{payoff}{payoff}:
	$$
	A^*_{\mathrm{Bayes}}
	=
	\arg\max_i \GTenToneHL{expected}{EMV}(\GTenToneHL{alternative}{A_i}).
	$$
	
	Therefore, the decision-maker selects the \GTenToneHL{alternative}{alternative} with the highest
	\GTenToneHL{expected}{Expected} \GTenToneHL{expected}{Monetary} \GTenToneHL{expected}{Value}.
	
	% ============================================================
	% GENERAL 2x2 PROBLEM
	% ============================================================
	
	\subsection{General 2x2 Problem}\label{sec:c4-intro-2x2}

	
	Consider the following \GTenToneHL{payoff}{payoff} table, where each \GTenToneHL{state}{state} \GTenToneHL{state}{of} \GTenToneHL{state}{nature} has an
	associated \GTenToneHL{probability}{probability}:
	
	\begin{center}
		\renewcommand{\arraystretch}{1.3}
		\begin{tabular}{lcc}
			\toprule
			Alternative & $\GTenToneHL{state}{S_1}$ & $\GTenToneHL{state}{S_2}$ \\
			\midrule
			$\GTenToneHL{alternative}{A_1}$
			& $\GTenToneHL{payoff}{P_{ay}}(\GTenToneHL{alternative}{A_1},\GTenToneHL{state}{S_1})$
			& $\GTenToneHL{payoff}{P_{ay}}(\GTenToneHL{alternative}{A_1},\GTenToneHL{state}{S_2})$ \\
			$\GTenToneHL{alternative}{A_2}$
			& $\GTenToneHL{payoff}{P_{ay}}(\GTenToneHL{alternative}{A_2},\GTenToneHL{state}{S_1})$
			& $\GTenToneHL{payoff}{P_{ay}}(\GTenToneHL{alternative}{A_2},\GTenToneHL{state}{S_2})$ \\
			\bottomrule
		\end{tabular}
		
	\end{center}
	
	The \GTenToneHL{probability}{probabilities} of the \GTenToneHL{state}{states} \GTenToneHL{state}{of} \GTenToneHL{state}{nature} are
	
	$$
	\GTenToneHL{probability}{P}(\GTenToneHL{state}{S_1}),
	\qquad
	\GTenToneHL{probability}{P}(\GTenToneHL{state}{S_2}),
	\qquad
	\GTenToneHL{probability}{P}(\GTenToneHL{state}{S_1})+\GTenToneHL{probability}{P}(\GTenToneHL{state}{S_2})=1.
	$$
	
	Determine the decision using the \GTenToneHL{expected}{Expected} \GTenToneHL{expected}{Monetary} \GTenToneHL{expected}{Value} criterion.
	
	% ============================================================
	% SOLUTION 2x2
	% ============================================================
	
	\subsubsection{Solution}
	
	\textbf{Step 1: Identify the \GTenToneHL{alternative}{alternatives}, \GTenToneHL{state}{states} \GTenToneHL{state}{of} \GTenToneHL{state}{nature}, and \GTenToneHL{probability}{probabilities}}
	
	Alternatives: $\GTenToneHL{alternative}{A_1}, \GTenToneHL{alternative}{A_2}$ \qquad
	States of nature: $\GTenToneHL{state}{S_1}, \GTenToneHL{state}{S_2}$ \qquad
	Probabilities
	$
	\GTenToneHL{probability}{P}(\GTenToneHL{state}{S_1}),
	\GTenToneHL{probability}{P}(\GTenToneHL{state}{S_2})
	$
	
	\noindent
	\setlength{\columnseprule}{0.4pt}
	
	\begin{multicols}{2}
		\vspace{0pt}
		
		
		\textbf{Step 2: Calculate the \GTenToneHL{expected}{EMV} of $\GTenToneHL{alternative}{A_1}$}
		
		For \GTenToneHL{alternative}{alternative} $\GTenToneHL{alternative}{A_1}$, multiply each \GTenToneHL{payoff}{payoff} by the \GTenToneHL{probability}{probability} of
		its corresponding \GTenToneHL{state}{state} \GTenToneHL{state}{of} \GTenToneHL{state}{nature} and add the results:
		
		\[
		\GTenToneHL{expected}{EMV}(\GTenToneHL{alternative}{A_1})
		=
		\GTenToneHL{probability}{P}(\GTenToneHL{state}{S_1})\GTenToneHL{payoff}{P_{ay}}(\GTenToneHL{alternative}{A_1},\GTenToneHL{state}{S_1})
		+
		\GTenToneHL{probability}{P}(\GTenToneHL{state}{S_2})\GTenToneHL{payoff}{P_{ay}}(\GTenToneHL{alternative}{A_1},\GTenToneHL{state}{S_2}).
		\]
		
		Thus, $\GTenToneHL{expected}{EMV}(\GTenToneHL{alternative}{A_1})$ represents the expected \GTenToneHL{payoff}{payoff} from selecting
		\GTenToneHL{alternative}{alternative} $\GTenToneHL{alternative}{A_1}$.
		
		\columnbreak
		
		\textbf{Step 3: Calculate the \GTenToneHL{expected}{EMV} of $\GTenToneHL{alternative}{A_2}$}
		
		For \GTenToneHL{alternative}{alternative} $\GTenToneHL{alternative}{A_2}$:
		
		\[
		\GTenToneHL{expected}{EMV}(\GTenToneHL{alternative}{A_2})
		=
		\GTenToneHL{probability}{P}(\GTenToneHL{state}{S_1})\GTenToneHL{payoff}{P_{ay}}(\GTenToneHL{alternative}{A_2},\GTenToneHL{state}{S_1})
		+
		\GTenToneHL{probability}{P}(\GTenToneHL{state}{S_2})\GTenToneHL{payoff}{P_{ay}}(\GTenToneHL{alternative}{A_2},\GTenToneHL{state}{S_2}).
		\]
		
		Thus, $\GTenToneHL{expected}{EMV}(\GTenToneHL{alternative}{A_2})$ represents the expected \GTenToneHL{payoff}{payoff} from selecting
		\GTenToneHL{alternative}{alternative} $\GTenToneHL{alternative}{A_2}$.
		
		
	\end{multicols}
	
	\setlength{\columnseprule}{0pt}
	
	\textbf{Step 4: Apply the Bayes decision criterion}
	
	Compare the two expected monetary values:
	
	$$
	A^*_{\mathrm{Bayes}}
	=
	\arg\max_i \GTenToneHL{expected}{EMV}(\GTenToneHL{alternative}{A_i}).
	$$
	
	Therefore, choose the \GTenToneHL{alternative}{alternative} with the highest Expected Monetary
	Value.
	
	% ============================================================
	% EXTENDED PAYOFF TABLE
	% ============================================================
	
	\textbf{Extended \GTenToneHL{payoff}{payoff} table}
	
	\begin{center}
		\small
		\renewcommand{\arraystretch}{1.5}
		\setlength{\tabcolsep}{4pt}
		
		
		\begin{tabular}{lccc}
			\toprule
			Alternative
			& $\GTenToneHL{state}{S_1}\;(\GTenToneHL{probability}{P}(\GTenToneHL{state}{S_1}))$
			& $\GTenToneHL{state}{S_2}\;(\GTenToneHL{probability}{P}(\GTenToneHL{state}{S_2}))$
			& \GTenToneHL{expected}{EMV}
			\\
			\midrule
			
			$\GTenToneHL{alternative}{A_1}$
			& $\GTenToneHL{payoff}{P_{ay}}(\GTenToneHL{alternative}{A_1},\GTenToneHL{state}{S_1})$
			& $\GTenToneHL{payoff}{P_{ay}}(\GTenToneHL{alternative}{A_1},\GTenToneHL{state}{S_2})$
			&
			$\GTenToneHL{probability}{P}(\GTenToneHL{state}{S_1})\GTenToneHL{payoff}{P_{ay}}(\GTenToneHL{alternative}{A_1},\GTenToneHL{state}{S_1})
			+
			\GTenToneHL{probability}{P}(\GTenToneHL{state}{S_2})\GTenToneHL{payoff}{P_{ay}}(\GTenToneHL{alternative}{A_1},\GTenToneHL{state}{S_2})$
			\\
			
			$\GTenToneHL{alternative}{A_2}$
			& $\GTenToneHL{payoff}{P_{ay}}(\GTenToneHL{alternative}{A_2},\GTenToneHL{state}{S_1})$
			& $\GTenToneHL{payoff}{P_{ay}}(\GTenToneHL{alternative}{A_2},\GTenToneHL{state}{S_2})$
			&
			$\GTenToneHL{probability}{P}(\GTenToneHL{state}{S_1})\GTenToneHL{payoff}{P_{ay}}(\GTenToneHL{alternative}{A_2},\GTenToneHL{state}{S_1})
			+
			\GTenToneHL{probability}{P}(\GTenToneHL{state}{S_2})\GTenToneHL{payoff}{P_{ay}}(\GTenToneHL{alternative}{A_2},\GTenToneHL{state}{S_2})$
			\\
			
			\bottomrule
		\end{tabular}
		
	\end{center}

	\subsection{General 4x3 Problem}\label{sec:c4-intro-4x3}
	
	Consider the following \GTenToneHL{payoff}{payoff} table, where each \GTenToneHL{state}{state} \GTenToneHL{state}{of} \GTenToneHL{state}{nature} has an
	associated \GTenToneHL{probability}{probability}:
	
	\begin{center}
		\renewcommand{\arraystretch}{1.3}
		\begin{tabular}{lccc}
			\toprule
			Alternative & $\GTenToneHL{state}{S_1}$ & $\GTenToneHL{state}{S_2}$ & $\GTenToneHL{state}{S_3}$ \\
			\midrule
		
			$\GTenToneHL{alternative}{A_1}$
			& $\GTenToneHL{payoff}{P_{ay}}(\GTenToneHL{alternative}{A_1},\GTenToneHL{state}{S_1})$
			& $\GTenToneHL{payoff}{P_{ay}}(\GTenToneHL{alternative}{A_1},\GTenToneHL{state}{S_2})$
			& $\GTenToneHL{payoff}{P_{ay}}(\GTenToneHL{alternative}{A_1},\GTenToneHL{state}{S_3})$
			\\
			
			$\GTenToneHL{alternative}{A_2}$
			& $\GTenToneHL{payoff}{P_{ay}}(\GTenToneHL{alternative}{A_2},\GTenToneHL{state}{S_1})$
			& $\GTenToneHL{payoff}{P_{ay}}(\GTenToneHL{alternative}{A_2},\GTenToneHL{state}{S_2})$
			& $\GTenToneHL{payoff}{P_{ay}}(\GTenToneHL{alternative}{A_2},\GTenToneHL{state}{S_3})$
			\\
			
			$\GTenToneHL{alternative}{A_3}$
			& $\GTenToneHL{payoff}{P_{ay}}(\GTenToneHL{alternative}{A_3},\GTenToneHL{state}{S_1})$
			& $\GTenToneHL{payoff}{P_{ay}}(\GTenToneHL{alternative}{A_3},\GTenToneHL{state}{S_2})$
			& $\GTenToneHL{payoff}{P_{ay}}(\GTenToneHL{alternative}{A_3},\GTenToneHL{state}{S_3})$
			\\
			
			$\GTenToneHL{alternative}{A_4}$
			& $\GTenToneHL{payoff}{P_{ay}}(\GTenToneHL{alternative}{A_4},\GTenToneHL{state}{S_1})$
			& $\GTenToneHL{payoff}{P_{ay}}(\GTenToneHL{alternative}{A_4},\GTenToneHL{state}{S_2})$
			& $\GTenToneHL{payoff}{P_{ay}}(\GTenToneHL{alternative}{A_4},\GTenToneHL{state}{S_3})$
			\\
			
			\bottomrule
		\end{tabular}
	\end{center}
	
	The \GTenToneHL{probability}{probabilities} of the \GTenToneHL{state}{states} \GTenToneHL{state}{of} \GTenToneHL{state}{nature} are
	
	$$
	\GTenToneHL{probability}{P}(\GTenToneHL{state}{S_1}), \qquad
	\GTenToneHL{probability}{P}(\GTenToneHL{state}{S_2}), \qquad
	\GTenToneHL{probability}{P}(\GTenToneHL{state}{S_3}),
	\qquad
	\text{with}
	\qquad 
	\GTenToneHL{probability}{P}(\GTenToneHL{state}{S_1})+\GTenToneHL{probability}{P}(\GTenToneHL{state}{S_2})+\GTenToneHL{probability}{P}(\GTenToneHL{state}{S_3})=1.
	$$
	
	Determine the decision using the \GTenToneHL{expected}{Expected} \GTenToneHL{expected}{Monetary} \GTenToneHL{expected}{Value} criterion.
	
	% ============================================================
	% SOLUTION 4x3
	% ============================================================
	
	\subsubsection{Solution}
	
	\textbf{Step 1: Identify the \GTenToneHL{alternative}{alternatives}, \GTenToneHL{state}{states} \GTenToneHL{state}{of} \GTenToneHL{state}{nature}, and \GTenToneHL{probability}{probabilities}}
	
	Alternatives: $\GTenToneHL{alternative}{A_1}, \GTenToneHL{alternative}{A_2}, \GTenToneHL{alternative}{A_3}, \GTenToneHL{alternative}{A_4}$
	\qquad
	States of nature: $\GTenToneHL{state}{S_1}, \GTenToneHL{state}{S_2}, \GTenToneHL{state}{S_3}$
	
	Probabilities:
	$
	\GTenToneHL{probability}{P}(\GTenToneHL{state}{S_1}), \GTenToneHL{probability}{P}(\GTenToneHL{state}{S_2}), \GTenToneHL{probability}{P}(\GTenToneHL{state}{S_3})
	$
	
	\noindent
	\setlength{\columnseprule}{0.4pt}
	
	\begin{multicols}{2}
		\vspace{0pt}
		
		\textbf{Step 2: Calculate the \GTenToneHL{expected}{EMV} of $\GTenToneHL{alternative}{A_1}$}
		\[
		\begin{aligned}
			\GTenToneHL{expected}{EMV}(\GTenToneHL{alternative}{A_1})
			&= \GTenToneHL{probability}{P}(\GTenToneHL{state}{S_1})\GTenToneHL{payoff}{P_{ay}}(\GTenToneHL{alternative}{A_1},\GTenToneHL{state}{S_1}) + \GTenToneHL{probability}{P}(\GTenToneHL{state}{S_2})\GTenToneHL{payoff}{P_{ay}}(\GTenToneHL{alternative}{A_1},\GTenToneHL{state}{S_2}) \\
			&\quad + \GTenToneHL{probability}{P}(\GTenToneHL{state}{S_3})\GTenToneHL{payoff}{P_{ay}}(\GTenToneHL{alternative}{A_1},\GTenToneHL{state}{S_3}).
		\end{aligned}
		\]
		
		\columnbreak
		
		\textbf{Step 3: Calculate the \GTenToneHL{expected}{EMV} of $\GTenToneHL{alternative}{A_2}$}
		\[
		\begin{aligned}
			\GTenToneHL{expected}{EMV}(\GTenToneHL{alternative}{A_2})
			&= \GTenToneHL{probability}{P}(\GTenToneHL{state}{S_1})\GTenToneHL{payoff}{P_{ay}}(\GTenToneHL{alternative}{A_2},\GTenToneHL{state}{S_1}) + \GTenToneHL{probability}{P}(\GTenToneHL{state}{S_2})\GTenToneHL{payoff}{P_{ay}}(\GTenToneHL{alternative}{A_2},\GTenToneHL{state}{S_2}) \\
			&\quad + \GTenToneHL{probability}{P}(\GTenToneHL{state}{S_3})\GTenToneHL{payoff}{P_{ay}}(\GTenToneHL{alternative}{A_2},\GTenToneHL{state}{S_3}).
		\end{aligned}
		\]
		
	\end{multicols}
	
	\begin{multicols}{2}
		\vspace{0pt}
		
		\textbf{Step 4: Calculate the \GTenToneHL{expected}{EMV} of $\GTenToneHL{alternative}{A_3}$}
		\[
		\begin{aligned}
			\GTenToneHL{expected}{EMV}(\GTenToneHL{alternative}{A_3})
			&= \GTenToneHL{probability}{P}(\GTenToneHL{state}{S_1})\GTenToneHL{payoff}{P_{ay}}(\GTenToneHL{alternative}{A_3},\GTenToneHL{state}{S_1}) + \GTenToneHL{probability}{P}(\GTenToneHL{state}{S_2})\GTenToneHL{payoff}{P_{ay}}(\GTenToneHL{alternative}{A_3},\GTenToneHL{state}{S_2}) \\
			&\quad + \GTenToneHL{probability}{P}(\GTenToneHL{state}{S_3})\GTenToneHL{payoff}{P_{ay}}(\GTenToneHL{alternative}{A_3},\GTenToneHL{state}{S_3}).
		\end{aligned}
		\]
		
		\columnbreak
		
		\textbf{Step 5: Calculate the \GTenToneHL{expected}{EMV} of $\GTenToneHL{alternative}{A_4}$}
		\[
		\begin{aligned}
			\GTenToneHL{expected}{EMV}(\GTenToneHL{alternative}{A_4})
			&= \GTenToneHL{probability}{P}(\GTenToneHL{state}{S_1})\GTenToneHL{payoff}{P_{ay}}(\GTenToneHL{alternative}{A_4},\GTenToneHL{state}{S_1}) + \GTenToneHL{probability}{P}(\GTenToneHL{state}{S_2})\GTenToneHL{payoff}{P_{ay}}(\GTenToneHL{alternative}{A_4},\GTenToneHL{state}{S_2}) \\
			&\quad + \GTenToneHL{probability}{P}(\GTenToneHL{state}{S_3})\GTenToneHL{payoff}{P_{ay}}(\GTenToneHL{alternative}{A_4},\GTenToneHL{state}{S_3}).
		\end{aligned}
		\]
		
	\end{multicols}
	
	\setlength{\columnseprule}{0pt}
	
	\textbf{Step 6: Apply the Bayes decision criterion}
	
	Compare the four expected monetary values:
	
	$$
	A^*_{\mathrm{Bayes}}
	=
	\arg\max_i \GTenToneHL{expected}{EMV}(\GTenToneHL{alternative}{A_i}).
	$$
	
	Therefore, choose the \GTenToneHL{alternative}{alternative} with the highest Expected Monetary
	Value.
	
	\textbf{Extended \GTenToneHL{payoff}{payoff} table}  \vspace{-3mm}
	\begin{center}
		\small
		\renewcommand{\arraystretch}{1.5}
		\setlength{\tabcolsep}{2pt}
		
		\begin{tabular}{|l|c|c|c|c|}
			\toprule
			Alternative
			& \begin{tabular}{c}$\GTenToneHL{state}{S_1}$ \\ $\GTenToneHL{probability}{P}(\GTenToneHL{state}{S_1})$\end{tabular}
			& \begin{tabular}{c}$\GTenToneHL{state}{S_2}$ \\ $\GTenToneHL{probability}{P}(\GTenToneHL{state}{S_2})$\end{tabular}
			& \begin{tabular}{c}$\GTenToneHL{state}{S_3}$ \\ $\GTenToneHL{probability}{P}(\GTenToneHL{state}{S_3})$\end{tabular}
			& \GTenToneHL{expected}{EMV}
			\\
			\midrule
			
			$\GTenToneHL{alternative}{A_1}$
			& $\GTenToneHL{payoff}{P_{ay}}(\GTenToneHL{alternative}{A_1},\GTenToneHL{state}{S_1})$
			& $\GTenToneHL{payoff}{P_{ay}}(\GTenToneHL{alternative}{A_1},\GTenToneHL{state}{S_2})$
			& $\GTenToneHL{payoff}{P_{ay}}(\GTenToneHL{alternative}{A_1},\GTenToneHL{state}{S_3})$
			&
			$\GTenToneHL{probability}{P}(\GTenToneHL{state}{S_1})\GTenToneHL{payoff}{P_{ay}}(\GTenToneHL{alternative}{A_1},\GTenToneHL{state}{S_1})
			+
			\GTenToneHL{probability}{P}(\GTenToneHL{state}{S_2})\GTenToneHL{payoff}{P_{ay}}(\GTenToneHL{alternative}{A_1},\GTenToneHL{state}{S_2})
			+
			\GTenToneHL{probability}{P}(\GTenToneHL{state}{S_3})\GTenToneHL{payoff}{P_{ay}}(\GTenToneHL{alternative}{A_1},\GTenToneHL{state}{S_3})$
			\\
			
			$\GTenToneHL{alternative}{A_2}$
			& $\GTenToneHL{payoff}{P_{ay}}(\GTenToneHL{alternative}{A_2},\GTenToneHL{state}{S_1})$
			& $\GTenToneHL{payoff}{P_{ay}}(\GTenToneHL{alternative}{A_2},\GTenToneHL{state}{S_2})$
			& $\GTenToneHL{payoff}{P_{ay}}(\GTenToneHL{alternative}{A_2},\GTenToneHL{state}{S_3})$
			&
			$\GTenToneHL{probability}{P}(\GTenToneHL{state}{S_1})\GTenToneHL{payoff}{P_{ay}}(\GTenToneHL{alternative}{A_2},\GTenToneHL{state}{S_1})
			+
			\GTenToneHL{probability}{P}(\GTenToneHL{state}{S_2})\GTenToneHL{payoff}{P_{ay}}(\GTenToneHL{alternative}{A_2},\GTenToneHL{state}{S_2})
			+
			\GTenToneHL{probability}{P}(\GTenToneHL{state}{S_3})\GTenToneHL{payoff}{P_{ay}}(\GTenToneHL{alternative}{A_2},\GTenToneHL{state}{S_3})$
			\\
			
			$\GTenToneHL{alternative}{A_3}$
			& $\GTenToneHL{payoff}{P_{ay}}(\GTenToneHL{alternative}{A_3},\GTenToneHL{state}{S_1})$
			& $\GTenToneHL{payoff}{P_{ay}}(\GTenToneHL{alternative}{A_3},\GTenToneHL{state}{S_2})$
			& $\GTenToneHL{payoff}{P_{ay}}(\GTenToneHL{alternative}{A_3},\GTenToneHL{state}{S_3})$
			&
			$\GTenToneHL{probability}{P}(\GTenToneHL{state}{S_1})\GTenToneHL{payoff}{P_{ay}}(\GTenToneHL{alternative}{A_3},\GTenToneHL{state}{S_1})
			+
			\GTenToneHL{probability}{P}(\GTenToneHL{state}{S_2})\GTenToneHL{payoff}{P_{ay}}(\GTenToneHL{alternative}{A_3},\GTenToneHL{state}{S_2})
			+
			\GTenToneHL{probability}{P}(\GTenToneHL{state}{S_3})\GTenToneHL{payoff}{P_{ay}}(\GTenToneHL{alternative}{A_3},\GTenToneHL{state}{S_3})$
			\\
			
			$\GTenToneHL{alternative}{A_4}$
			& $\GTenToneHL{payoff}{P_{ay}}(\GTenToneHL{alternative}{A_4},\GTenToneHL{state}{S_1})$
			& $\GTenToneHL{payoff}{P_{ay}}(\GTenToneHL{alternative}{A_4},\GTenToneHL{state}{S_2})$
			& $\GTenToneHL{payoff}{P_{ay}}(\GTenToneHL{alternative}{A_4},\GTenToneHL{state}{S_3})$
			&
			$\GTenToneHL{probability}{P}(\GTenToneHL{state}{S_1})\GTenToneHL{payoff}{P_{ay}}(\GTenToneHL{alternative}{A_4},\GTenToneHL{state}{S_1})
			+
			\GTenToneHL{probability}{P}(\GTenToneHL{state}{S_2})\GTenToneHL{payoff}{P_{ay}}(\GTenToneHL{alternative}{A_4},\GTenToneHL{state}{S_2})
			+
			\GTenToneHL{probability}{P}(\GTenToneHL{state}{S_3})\GTenToneHL{payoff}{P_{ay}}(\GTenToneHL{alternative}{A_4},\GTenToneHL{state}{S_3})$
			\\
			
			\bottomrule
		\end{tabular}
	\end{center}
	
	
\phantomsection\label{tone:reference}
\subsection*{Notation reference}
\GTenToneReference{C6}
\section{Contextualized Problems}\label{sec:c4-introduction}
	\SectionProblemsTableOfContents{
		\SectionProblemLink{sec:c4-intro-coffee}{2x2 Problem - Coffee shop} \\
		\SectionProblemLink{sec:c4-intro-clothing}{2x2 Problem - Clothing store} \\
		\SectionProblemLink{sec:c4-intro-technology}{4x3 Problem - Technology company} \\
		\SectionProblemLink{sec:c4-intro-logistics}{3x4 Problem - Logistics company}
	}\vspace{8mm}

	\subsection{2x2 Problem - Coffee shop}\label{sec:c4-intro-coffee}

	
	A small coffee shop is considering whether to introduce a new dessert line
	or keep its current menu. The manager estimates that there is a $60\%$
	\GTenToneHL{probability}{probability} of high customer traffic during the promotional weekend and a
	$40\%$ \GTenToneHL{probability}{probability} of low customer traffic. All \GTenToneHL{payoff}{payoffs} are expressed in thousands of dollars. What would a perfectly logical decision-maker choose using the Expected
	Monetary Value criterion? \vspace{-2mm}
	\begin{center}
		\renewcommand{\arraystretch}{1.3}
		\begin{tabular}{lcc}
			\toprule
			Alternative
			& \begin{tabular}{c}High traffic $\GTenToneHL{state}{S_1}$ \\ $\GTenToneHL{probability}{P}(\GTenToneHL{state}{S_1})=\GTenToneHL{probability}{0.60}$\end{tabular}
			& \begin{tabular}{c}Low traffic $\GTenToneHL{state}{S_2}$ \\ $\GTenToneHL{probability}{P}(\GTenToneHL{state}{S_2})=\GTenToneHL{probability}{0.40}$\end{tabular}
			\\
			\midrule
			Launch dessert line ($\GTenToneHL{alternative}{A_1}$) & 42 & 8 \\
			Keep current menu ($\GTenToneHL{alternative}{A_2}$) & 28 & 20 \\
			\bottomrule
		\end{tabular}
	\end{center} \vspace{-5mm}
	
	% ============================================================
	% SOLUTION 2x2
	% ============================================================
	
	\subsubsection{Solution}
	
	\textbf{Step 1: Identify the \GTenToneHL{alternative}{alternatives}, \GTenToneHL{state}{states} \GTenToneHL{state}{of} \GTenToneHL{state}{nature}, and \GTenToneHL{probability}{probabilities}}
	
	Alternatives:
	$\GTenToneHL{alternative}{A_1}$ = Launch dessert line,
	$\GTenToneHL{alternative}{A_2}$ = Keep current menu
	
	States of nature:
	$\GTenToneHL{state}{S_1}$ = High traffic,
	$\GTenToneHL{state}{S_2}$ = Low traffic
	
	Probabilities:
	$
	\GTenToneHL{probability}{P}(\GTenToneHL{state}{S_1})=\GTenToneHL{probability}{0.60},
	\GTenToneHL{probability}{P}(\GTenToneHL{state}{S_2})=\GTenToneHL{probability}{0.40}
	$ \vspace{-3mm}
	
	\noindent
	\setlength{\columnseprule}{0.4pt} 
	\begin{multicols}{2}
		\vspace{0pt}
		
		\textbf{Step 2: Calculate the \GTenToneHL{expected}{EMV} of $\GTenToneHL{alternative}{A_1}$}
		
		For the Launch dessert line \GTenToneHL{alternative}{alternative}:
		\[
		\GTenToneHL{expected}{EMV}(\GTenToneHL{alternative}{A_1})
		=
		\GTenToneHL{probability}{P}(\GTenToneHL{state}{S_1})(42)
		+
		\GTenToneHL{probability}{P}(\GTenToneHL{state}{S_2})(8)
		\]
		\[
		=
		(0.60)(42)
		+
		(0.40)(8)
		\]
		\[
		=
		25.2+3.2
		=
		28.4.
		\]
		
		Thus, the expected \GTenToneHL{payoff}{payoff} from launching the dessert line is
		$28.4$ thousand dollars.
		
		\columnbreak
		
		\textbf{Step 3: Calculate the \GTenToneHL{expected}{EMV} of $\GTenToneHL{alternative}{A_2}$}
		
		For the Keep current menu \GTenToneHL{alternative}{alternative}:
		\[
		\GTenToneHL{expected}{EMV}(\GTenToneHL{alternative}{A_2})
		=
		\GTenToneHL{probability}{P}(\GTenToneHL{state}{S_1})(28)
		+
		\GTenToneHL{probability}{P}(\GTenToneHL{state}{S_2})(20)
		\]
		\[
		=
		(0.60)(28)
		+
		(0.40)(20)
		\]
		\[
		=
		16.8+8
		=
		24.8.
		\]
		
		Thus, the expected \GTenToneHL{payoff}{payoff} from keeping the current menu is
		$24.8$ thousand dollars.
		
	\end{multicols}
	
	\setlength{\columnseprule}{0pt}
	
	\newpage
	
	\textbf{Step 4: Apply the Bayes decision criterion}
	
	Compare the two expected monetary values:
	$$
	\GTenToneHL{expected}{EMV}(\GTenToneHL{alternative}{A_1})=\GTenToneHL{expected}{28.4},
	\qquad
	\GTenToneHL{expected}{EMV}(\GTenToneHL{alternative}{A_2})=24.8.
	$$
	$$
	A^*_{\mathrm{Bayes}}
	=
	\arg\max_i \GTenToneHL{expected}{EMV}(\GTenToneHL{alternative}{A_i})
	=
	\GTenToneHL{alternative}{A_1}.
	$$
	
	Thus, the Bayes criterion recommends launching the dessert line because
	it has the highest \GTenToneHL{expected}{Expected} \GTenToneHL{expected}{Monetary} \GTenToneHL{expected}{Value}, with an expected \GTenToneHL{payoff}{payoff} of
	$28.4$ thousand dollars.
	
	% ============================================================
	% EXTENDED PAYOFF TABLE
	% ============================================================
	
	\textbf{Extended \GTenToneHL{payoff}{payoff} table}
	\textbf{Extended \GTenToneHL{payoff}{payoff} table}
	
	\begin{center}
		\renewcommand{\arraystretch}{1.5}
		\setlength{\tabcolsep}{4pt}
		
		\begin{tabular}{lccc}
			\toprule
			Alternative
			& \begin{tabular}{c}High traffic $\GTenToneHL{state}{S_1}$ \\ $\GTenToneHL{probability}{P}(\GTenToneHL{state}{S_1})=\GTenToneHL{probability}{0.60}$\end{tabular}
			& \begin{tabular}{c}Low traffic $\GTenToneHL{state}{S_2}$ \\ $\GTenToneHL{probability}{P}(\GTenToneHL{state}{S_2})=\GTenToneHL{probability}{0.40}$\end{tabular}
			& \GTenToneHL{expected}{EMV}
			\\
			\midrule
			
			$\GTenToneHL{alternative}{A_1}$: Launch dessert line
			& $42$ & $8$
			& $0.60(42)+0.40(8)=28.4$
			\\
			
			$\GTenToneHL{alternative}{A_2}$: Keep current menu
			& $28$ & $20$
			& $0.60(28)+0.40(20)=24.8$
			\\
			
			\bottomrule
		\end{tabular}
	\end{center}
	
	\subsection{2x2 Problem - Clothing store}\label{sec:c4-intro-clothing}

	
	A local clothing store is deciding whether to launch a new advertising campaign
	before the holiday season or continue with its current marketing strategy. The
	manager estimates that there is a $70\%$ \GTenToneHL{probability}{probability} of strong sales during
	the holiday season and a $30\%$ \GTenToneHL{probability}{probability} of weak sales. All \GTenToneHL{payoff}{payoffs} are
	expressed in thousands of dollars.
	
	What would a perfectly logical decision-maker choose using the Expected
	Monetary Value criterion?
	
	\begin{center}
		\renewcommand{\arraystretch}{1.3}
		\begin{tabular}{lcc}
			\toprule
			Alternative
			& \begin{tabular}{c}
				Strong sales $\GTenToneHL{state}{S_1}$ \\
				$\GTenToneHL{probability}{P}(\GTenToneHL{state}{S_1})=\GTenToneHL{probability}{0.70}$
			\end{tabular}
			& \begin{tabular}{c}
				Weak sales $\GTenToneHL{state}{S_2}$ \\
				$\GTenToneHL{probability}{P}(\GTenToneHL{state}{S_2})=\GTenToneHL{probability}{0.30}$
			\end{tabular}
			\\
			\midrule
			Launch new campaign ($\GTenToneHL{alternative}{A_1}$) & 55 & 5 \\
			Keep current strategy ($\GTenToneHL{alternative}{A_2}$) & 35 & 22 \\
			\bottomrule
		\end{tabular}
	\end{center}
	
	% ============================================================
	% SOLUTION 2x2
	% ============================================================
	
	\subsubsection{Solution}
	
	\textbf{Step 1: Identify the \GTenToneHL{alternative}{alternatives}, \GTenToneHL{state}{states} \GTenToneHL{state}{of} \GTenToneHL{state}{nature}, and \GTenToneHL{probability}{probabilities}}
	
	Alternatives:
	$\GTenToneHL{alternative}{A_1}$ = Launch new campaign,
	$\GTenToneHL{alternative}{A_2}$ = Keep current strategy
	
	States of nature:
	$\GTenToneHL{state}{S_1}$ = Strong sales,
	$\GTenToneHL{state}{S_2}$ = Weak sales
	
	Probabilities:
	$
	\GTenToneHL{probability}{P}(\GTenToneHL{state}{S_1})=\GTenToneHL{probability}{0.70},
	\GTenToneHL{probability}{P}(\GTenToneHL{state}{S_2})=\GTenToneHL{probability}{0.30}
	$
	
	\noindent
	\setlength{\columnseprule}{0.4pt}
	\newpage
	
	\begin{multicols}{2}
		\vspace{0pt}
		
		\textbf{Step 2: Calculate the \GTenToneHL{expected}{EMV} of $\GTenToneHL{alternative}{A_1}$}
		
		For the Launch new campaign \GTenToneHL{alternative}{alternative}:
		\[
		\GTenToneHL{expected}{EMV}(\GTenToneHL{alternative}{A_1})
		=
		\GTenToneHL{probability}{P}(\GTenToneHL{state}{S_1})(55)
		+
		\GTenToneHL{probability}{P}(\GTenToneHL{state}{S_2})(5)
		\]
		\[
		=
		(0.70)(55)
		+
		(0.30)(5)
		\]
		\[
		=
		38.5+1.5
		=
		40.
		\]
		
		Thus, the expected \GTenToneHL{payoff}{payoff} from launching the new campaign is
		$40$ thousand dollars.
		
		\columnbreak
		
		\textbf{Step 3: Calculate the \GTenToneHL{expected}{EMV} of $\GTenToneHL{alternative}{A_2}$}
		
		For the Keep current strategy \GTenToneHL{alternative}{alternative}:
		\[
		\GTenToneHL{expected}{EMV}(\GTenToneHL{alternative}{A_2})
		=
		\GTenToneHL{probability}{P}(\GTenToneHL{state}{S_1})(35)
		+
		\GTenToneHL{probability}{P}(\GTenToneHL{state}{S_2})(22)
		\]
		\[
		=
		(0.70)(35)
		+
		(0.30)(22)
		\]
		\[
		=
		24.5+6.6
		=
		31.1.
		\]
		
		Thus, the expected \GTenToneHL{payoff}{payoff} from keeping the current strategy is
		$31.1$ thousand dollars.
		
	\end{multicols}
	
	\setlength{\columnseprule}{0pt}
	
	\textbf{Step 4: Apply the Bayes decision criterion}
	
	Compare the two expected monetary values:
	$$
	\GTenToneHL{expected}{EMV}(\GTenToneHL{alternative}{A_1})=\GTenToneHL{expected}{40},
	\qquad
	\GTenToneHL{expected}{EMV}(\GTenToneHL{alternative}{A_2})=31.1.
	$$
	$$
	A^*_{\mathrm{Bayes}}
	=
	\arg\max_i \GTenToneHL{expected}{EMV}(\GTenToneHL{alternative}{A_i})
	=
	\GTenToneHL{alternative}{A_1}.
	$$
	
	Thus, the Bayes criterion recommends launching the new campaign because
	it has the highest \GTenToneHL{expected}{Expected} \GTenToneHL{expected}{Monetary} \GTenToneHL{expected}{Value}, with an expected \GTenToneHL{payoff}{payoff} of
	$40$ thousand dollars.
	
	% ============================================================
	% EXTENDED PAYOFF TABLE
	% ============================================================
	
	\textbf{Extended \GTenToneHL{payoff}{payoff} table}
	
	\begin{center}
		\renewcommand{\arraystretch}{1.5}
		\setlength{\tabcolsep}{4pt}
		
		\begin{tabular}{lccc}
			\toprule
			Alternative
			& \begin{tabular}{c}$\GTenToneHL{state}{S_1}$ \\ $\GTenToneHL{probability}{P}(\GTenToneHL{state}{S_1})=\GTenToneHL{probability}{0.70}$\end{tabular}
			& \begin{tabular}{c}$\GTenToneHL{state}{S_2}$ \\ $\GTenToneHL{probability}{P}(\GTenToneHL{state}{S_2})=\GTenToneHL{probability}{0.30}$\end{tabular}
			& \GTenToneHL{expected}{EMV}
			\\
			\midrule
			
			$\GTenToneHL{alternative}{A_1}$: Launch new campaign
			& $55$
			& $5$
			& $0.70(55)+0.30(5)=40$
			\\
			
			$\GTenToneHL{alternative}{A_2}$: Keep current strategy
			& $35$
			& $22$
			& $0.70(35)+0.30(22)=31.1$
			\\
			
			\bottomrule
		\end{tabular}
	\end{center}

	\subsection{4x3 Problem - Technology company}\label{sec:c4-intro-technology}

	
	A technology company is deciding which product strategy to adopt for the launch
	of a new device. The manager estimates that there is a $50\%$ \GTenToneHL{probability}{probability} of
	high demand, a $30\%$ \GTenToneHL{probability}{probability} of medium demand, and a $20\%$ \GTenToneHL{probability}{probability}
	of low demand. The company is considering four possible strategies: launching
	a premium version, launching a standard version, launching a budget version,
	or delaying the launch.	The estimated \GTenToneHL{payoff}{profit} for each strategy under each possible market condition is
	shown in the table below. All \GTenToneHL{payoff}{payoffs} are expressed in thousands of dollars.
	
	What would a perfectly logical decision-maker choose using the Expected
	Monetary Value criterion?
	
	\begin{center}
		\renewcommand{\arraystretch}{1.3}
		\begin{tabular}{lccc}
			\toprule
			Alternative
			& \begin{tabular}{c}High demand $\GTenToneHL{state}{S_1}$ \\ $\GTenToneHL{probability}{P}(\GTenToneHL{state}{S_1})=\GTenToneHL{probability}{0.50}$\end{tabular}
			& \begin{tabular}{c}Medium demand $\GTenToneHL{state}{S_2}$ \\ $\GTenToneHL{probability}{P}(\GTenToneHL{state}{S_2})=\GTenToneHL{probability}{0.30}$\end{tabular}
			& \begin{tabular}{c}Low demand $\GTenToneHL{state}{S_3}$ \\ $\GTenToneHL{probability}{P}(\GTenToneHL{state}{S_3})=\GTenToneHL{probability}{0.20}$\end{tabular}
			\\
			\midrule
			Premium version ($\GTenToneHL{alternative}{A_1}$) & 80 & 35 & -20 \\
			Standard version ($\GTenToneHL{alternative}{A_2}$) & 60 & 42 & 18 \\
			Budget version ($\GTenToneHL{alternative}{A_3}$) & 40 & 32 & 25 \\
			Delay launch ($\GTenToneHL{alternative}{A_4}$) & 15 & 15 & 15 \\
			\bottomrule
		\end{tabular}
	\end{center}
	
	
	
	% ============================================================
	% SOLUTION 4x3
	% ============================================================
	
	\subsubsection{Solution}
	
	\textbf{Step 1: Identify the \GTenToneHL{alternative}{alternatives}, \GTenToneHL{state}{states} \GTenToneHL{state}{of} \GTenToneHL{state}{nature}, and \GTenToneHL{probability}{probabilities}}
	
	Alternatives:
	$\GTenToneHL{alternative}{A_1}$ = Premium version,
	$\GTenToneHL{alternative}{A_2}$ = Standard version,
	$\GTenToneHL{alternative}{A_3}$ = Budget version,
	$\GTenToneHL{alternative}{A_4}$ = Delay launch
	
	States of nature:
	$\GTenToneHL{state}{S_1}$ = High demand,
	$\GTenToneHL{state}{S_2}$ = Medium demand,
	$\GTenToneHL{state}{S_3}$ = Low demand
	
	Probabilities:
	$
	\GTenToneHL{probability}{P}(\GTenToneHL{state}{S_1})=\GTenToneHL{probability}{0.50},
	\GTenToneHL{probability}{P}(\GTenToneHL{state}{S_2})=\GTenToneHL{probability}{0.30},
	\GTenToneHL{probability}{P}(\GTenToneHL{state}{S_3})=\GTenToneHL{probability}{0.20}
	$
	
	\noindent
	\setlength{\columnseprule}{0.4pt}
	
	\begin{multicols}{2}
		\vspace{0pt}
		
		\textbf{Step 2: Calculate the \GTenToneHL{expected}{EMV} of $\GTenToneHL{alternative}{A_1}$ and $\GTenToneHL{alternative}{A_2}$}
		
		For the Premium version \GTenToneHL{alternative}{alternative}:
		\[
		\GTenToneHL{expected}{EMV}(\GTenToneHL{alternative}{A_1})
		=
		\GTenToneHL{probability}{P}(\GTenToneHL{state}{S_1})(80)
		+
		\GTenToneHL{probability}{P}(\GTenToneHL{state}{S_2})(35)
		+
		\GTenToneHL{probability}{P}(\GTenToneHL{state}{S_3})(-20)
		\]
		\[
		=
		(0.50)(80)
		+
		(0.30)(35)
		+
		(0.20)(-20)
		\]
		\[
		=
		40+10.5-4
		=
		46.5
		\]
		
		\vspace{2mm}
		
		For the Standard version \GTenToneHL{alternative}{alternative}:
		\[
		\GTenToneHL{expected}{EMV}(\GTenToneHL{alternative}{A_2})
		=
		\GTenToneHL{probability}{P}(\GTenToneHL{state}{S_1})(60)
		+
		\GTenToneHL{probability}{P}(\GTenToneHL{state}{S_2})(42)
		+
		\GTenToneHL{probability}{P}(\GTenToneHL{state}{S_3})(18)
		\]
		\[
		=
		(0.50)(60)
		+
		(0.30)(42)
		+
		(0.20)(18)
		\]
		\[
		=
		30+12.6+3.6
		=
		46.2
		\]
		
		\columnbreak
		
		\textbf{Step 3: Calculate the \GTenToneHL{expected}{EMV} of $\GTenToneHL{alternative}{A_3}$ and $\GTenToneHL{alternative}{A_4}$}
		
		For the Budget version \GTenToneHL{alternative}{alternative}:
		\[
		\GTenToneHL{expected}{EMV}(\GTenToneHL{alternative}{A_3})
		=
		\GTenToneHL{probability}{P}(\GTenToneHL{state}{S_1})(40)
		+
		\GTenToneHL{probability}{P}(\GTenToneHL{state}{S_2})(32)
		+
		\GTenToneHL{probability}{P}(\GTenToneHL{state}{S_3})(25)
		\]
		\[
		=
		(0.50)(40)
		+
		(0.30)(32)
		+
		(0.20)(25)
		\]
		\[
		=
		20+9.6+5
		=
		34.6
		\]
		
		\vspace{2mm}
		
		For the Delay launch \GTenToneHL{alternative}{alternative}:
		\[
		\GTenToneHL{expected}{EMV}(\GTenToneHL{alternative}{A_4})
		=
		\GTenToneHL{probability}{P}(\GTenToneHL{state}{S_1})(15)
		+
		\GTenToneHL{probability}{P}(\GTenToneHL{state}{S_2})(15)
		+
		\GTenToneHL{probability}{P}(\GTenToneHL{state}{S_3})(15)
		\]
		\[
		=
		(0.50)(15)
		+
		(0.30)(15)
		+
		(0.20)(15)
		\]
		\[
		=
		7.5+4.5+3
		=
		15
		\]
		
	\end{multicols}
	
	\setlength{\columnseprule}{0pt}
	
	\textbf{Step 4: Apply the Bayes decision criterion}
	
	Compare the four expected monetary values:
	$$
	\GTenToneHL{expected}{EMV}(\GTenToneHL{alternative}{A_1})=\GTenToneHL{expected}{46.5}
	\qquad
	\GTenToneHL{expected}{EMV}(\GTenToneHL{alternative}{A_2})=\GTenToneHL{expected}{46.2}
	\qquad
	\GTenToneHL{expected}{EMV}(\GTenToneHL{alternative}{A_3})=\GTenToneHL{expected}{34.6}
	\qquad
	\GTenToneHL{expected}{EMV}(\GTenToneHL{alternative}{A_4})=\GTenToneHL{expected}{15}
	$$
	$$
	A^*_{\mathrm{Bayes}}
	=
	\arg\max_i \GTenToneHL{expected}{EMV}(\GTenToneHL{alternative}{A_i})
	=
	\GTenToneHL{alternative}{A_1}
	$$
	
	Thus, the Bayes criterion recommends launching the premium version because
	it has the highest \GTenToneHL{expected}{Expected} \GTenToneHL{expected}{Monetary} \GTenToneHL{expected}{Value}, with an expected \GTenToneHL{payoff}{payoff} of
	$46.5$ thousand dollars.
	
	% ============================================================
	% EXTENDED PAYOFF TABLE
	% ============================================================
	
	\newpage
	\textbf{Extended \GTenToneHL{payoff}{payoff} table}
	
	\begin{center}
		\normalsize
		\renewcommand{\arraystretch}{1.5}
		\setlength{\tabcolsep}{3pt}
		
		\begin{tabular}{lcccc}
			\toprule
			Alternative
			& \begin{tabular}{c}$\GTenToneHL{state}{S_1}$ \\ $\GTenToneHL{probability}{P}(\GTenToneHL{state}{S_1})=\GTenToneHL{probability}{0.5}$\end{tabular}
			& \begin{tabular}{c}$\GTenToneHL{state}{S_2}$ \\ $\GTenToneHL{probability}{P}(\GTenToneHL{state}{S_2})=\GTenToneHL{probability}{0.3}$\end{tabular}
			& \begin{tabular}{c}$\GTenToneHL{state}{S_3}$ \\ $\GTenToneHL{probability}{P}(\GTenToneHL{state}{S_3})=\GTenToneHL{probability}{0.2}$\end{tabular}
			& \GTenToneHL{expected}{EMV}
			\\
			\midrule
			
			$\GTenToneHL{alternative}{A_1}$: Premium version
			& $80$ & $35$ & $-20$
			& $0.5(80)+0.3(35)+0.2(-20)=46.5$
			\\
			
			$\GTenToneHL{alternative}{A_2}$: Standard version
			& $60$ & $42$ & $18$
			& $0.5(60)+0.3(42)+0.2(18)=46.2$
			\\
			
			$\GTenToneHL{alternative}{A_3}$: Budget version
			& $40$ & $32$ & $25$
			& $0.5(40)+0.3(32)+0.2(25)=34.6$
			\\
			
			$\GTenToneHL{alternative}{A_4}$: Delay launch
			& $15$ & $15$ & $15$
			& $0.5(15)+0.3(15)+0.2(15)=15$
			\\
			
			\bottomrule
		\end{tabular}
		\end{center}
	
	\subsection{3x4 Problem - Logistics company}\label{sec:c4-intro-logistics}
	
	A logistics company is planning its warehouse capacity for the next year.
	The manager expects demand for storage services to vary depending on the
	level of activity in the market. Based on current forecasts, there is a
	$20\%$ \GTenToneHL{probability}{probability} of very low demand, a $30\%$ \GTenToneHL{probability}{probability} of low demand,
	a $35\%$ \GTenToneHL{probability}{probability} of high demand, and a $15\%$ \GTenToneHL{probability}{probability} of very high
	demand.
	
	The company is considering three possible capacity strategies: renting a
	small warehouse, renting a medium warehouse, or renting a large warehouse.
	The estimated annual \GTenToneHL{payoff}{profit} for each strategy under each possible demand
	level is shown in the table below. All \GTenToneHL{payoff}{payoffs} are expressed in thousands
	of dollars.
	
	What would a perfectly logical decision-maker choose using the Expected
	Monetary Value criterion?
	
	\begin{center}
		\renewcommand{\arraystretch}{1.3}
		\begin{tabular}{lcccc}
			\toprule
			Alternative
			& \begin{tabular}{c}Very low demand $\GTenToneHL{state}{S_1}$ \\ $\GTenToneHL{probability}{P}(\GTenToneHL{state}{S_1})=\GTenToneHL{probability}{0.20}$\end{tabular}
			& \begin{tabular}{c}Low demand $\GTenToneHL{state}{S_2}$ \\ $\GTenToneHL{probability}{P}(\GTenToneHL{state}{S_2})=\GTenToneHL{probability}{0.30}$\end{tabular}
			& \begin{tabular}{c}High demand $\GTenToneHL{state}{S_3}$ \\ $\GTenToneHL{probability}{P}(\GTenToneHL{state}{S_3})=\GTenToneHL{probability}{0.35}$\end{tabular}
			& \begin{tabular}{c}Very high demand $\GTenToneHL{state}{S_4}$ \\ $\GTenToneHL{probability}{P}(\GTenToneHL{state}{S_4})=\GTenToneHL{probability}{0.15}$\end{tabular}
			\\
			\midrule
			Small warehouse ($\GTenToneHL{alternative}{A_1}$) & 25 & 35 & 48 & 55 \\
			Medium warehouse ($\GTenToneHL{alternative}{A_2}$) & 5 & 32 & 65 & 82 \\
			Large warehouse ($\GTenToneHL{alternative}{A_3}$) & -25 & 15 & 72 & 110 \\
			\bottomrule
		\end{tabular}
	\end{center}
	
	
	% ============================================================
	% SOLUTION 3x4
	% ============================================================
	
	\subsubsection{Solution}
	
	\textbf{Step 1: Identify the \GTenToneHL{alternative}{alternatives}, \GTenToneHL{state}{states} \GTenToneHL{state}{of} \GTenToneHL{state}{nature}, and \GTenToneHL{probability}{probabilities}}
	
	Alternatives:
	$\GTenToneHL{alternative}{A_1}$ = Small warehouse,
	$\GTenToneHL{alternative}{A_2}$ = Medium warehouse,
	$\GTenToneHL{alternative}{A_3}$ = Large warehouse
	
	States of nature:
	$\GTenToneHL{state}{S_1}$ = Very low demand,
	$\GTenToneHL{state}{S_2}$ = Low demand,
	$\GTenToneHL{state}{S_3}$ = High demand,
	$\GTenToneHL{state}{S_4}$ = Very high demand
	
	Probabilities:
	$
	\GTenToneHL{probability}{P}(\GTenToneHL{state}{S_1})=\GTenToneHL{probability}{0.20},
	\GTenToneHL{probability}{P}(\GTenToneHL{state}{S_2})=\GTenToneHL{probability}{0.30},
	\GTenToneHL{probability}{P}(\GTenToneHL{state}{S_3})=\GTenToneHL{probability}{0.35},
	\GTenToneHL{probability}{P}(\GTenToneHL{state}{S_4})=\GTenToneHL{probability}{0.15}
	$
	
	\noindent
	\setlength{\columnseprule}{0.4pt}
	
	\begin{multicols}{2}
		\vspace{0pt}
		
		\textbf{Step 2: Calculate the \GTenToneHL{expected}{EMV} of $\GTenToneHL{alternative}{A_1}$ and $\GTenToneHL{alternative}{A_2}$}
		
		For the Small warehouse \GTenToneHL{alternative}{alternative}:
		\[
		\GTenToneHL{expected}{EMV}(\GTenToneHL{alternative}{A_1})
		=
		\GTenToneHL{probability}{P}(\GTenToneHL{state}{S_1})(25)
		+
		\GTenToneHL{probability}{P}(\GTenToneHL{state}{S_2})(35)
		+
		\GTenToneHL{probability}{P}(\GTenToneHL{state}{S_3})(48)
		+
		\GTenToneHL{probability}{P}(\GTenToneHL{state}{S_4})(55)
		\]
		\[
		=
		(0.20)(25)
		+
		(0.30)(35)
		+
		(0.35)(48)
		+
		(0.15)(55)
		\]
		\[
		=
		5+10.5+16.8+8.25
		=
		40.55
		\]
		
		\vspace{2mm}
		
		For the Medium warehouse \GTenToneHL{alternative}{alternative}:
		\[
		\GTenToneHL{expected}{EMV}(\GTenToneHL{alternative}{A_2})
		=
		\GTenToneHL{probability}{P}(\GTenToneHL{state}{S_1})(5)
		+
		\GTenToneHL{probability}{P}(\GTenToneHL{state}{S_2})(32)
		+
		\GTenToneHL{probability}{P}(\GTenToneHL{state}{S_3})(65)
		+
		\GTenToneHL{probability}{P}(\GTenToneHL{state}{S_4})(82)
		\]
		\[
		=
		(0.20)(5)
		+
		(0.30)(32)
		+
		(0.35)(65)
		+
		(0.15)(82)
		\]
		\[
		=
		1+9.6+22.75+12.3
		=
		45.65
		\]
		
		\columnbreak
		
		\textbf{Step 3: Calculate the \GTenToneHL{expected}{EMV} of $\GTenToneHL{alternative}{A_3}$}
		
		For the Large warehouse \GTenToneHL{alternative}{alternative}:
		\[
		\GTenToneHL{expected}{EMV}(\GTenToneHL{alternative}{A_3})
		=
		\GTenToneHL{probability}{P}(\GTenToneHL{state}{S_1})(-25)
		+
		\GTenToneHL{probability}{P}(\GTenToneHL{state}{S_2})(15)
		+
		\GTenToneHL{probability}{P}(\GTenToneHL{state}{S_3})(72)
		+
		\GTenToneHL{probability}{P}(\GTenToneHL{state}{S_4})(110)
		\]
		\[
		=
		(0.20)(-25)
		+
		(0.30)(15)
		+
		(0.35)(72)
		+
		(0.15)(110)
		\]
		\[
		=
		-5+4.5+25.2+16.5
		=
		41.2
		\]
		
	\end{multicols}
	
	\setlength{\columnseprule}{0pt}
	
	\textbf{Step 4: Apply the Bayes decision criterion}
	
	Compare the three expected monetary values:
	$$
	\GTenToneHL{expected}{EMV}(\GTenToneHL{alternative}{A_1})=\GTenToneHL{expected}{40.55}
	\qquad
	\GTenToneHL{expected}{EMV}(\GTenToneHL{alternative}{A_2})=\GTenToneHL{expected}{45.65}
	\qquad
	\GTenToneHL{expected}{EMV}(\GTenToneHL{alternative}{A_3})=\GTenToneHL{expected}{41.2}
	$$
	$$
	A^*_{\mathrm{Bayes}}
	=
	\arg\max_i \GTenToneHL{expected}{EMV}(\GTenToneHL{alternative}{A_i})
	=
	\GTenToneHL{alternative}{A_2}
	$$
	
	Thus, the Bayes criterion recommends renting the medium warehouse because
	it has the highest \GTenToneHL{expected}{Expected} \GTenToneHL{expected}{Monetary} \GTenToneHL{expected}{Value}, with an expected annual \GTenToneHL{payoff}{payoff}
	of $45.65$ thousand dollars.
	
	% ============================================================
	% EXTENDED PAYOFF TABLE
	% ============================================================
	
	\textbf{Extended \GTenToneHL{payoff}{payoff} table}
	
	\begin{center}
		\normalsize
		\renewcommand{\arraystretch}{1.5}
		\setlength{\tabcolsep}{3pt}
		
		\begin{tabular}{l|c|c|c|c}
			\toprule
			State
			& $\GTenToneHL{alternative}{A_1}$: Small warehouse
			& $\GTenToneHL{alternative}{A_2}$: Medium warehouse
			& $\GTenToneHL{alternative}{A_3}$: Large warehouse
			& Probability
			\\
			\midrule
			
			$\GTenToneHL{state}{S_1}$
			& $25$
			& $5$
			& $-25$
			& $0.2$
			\\
			
			$\GTenToneHL{state}{S_2}$
			& $35$
			& $32$
			& $15$
			& $0.3$
			\\
			
			$\GTenToneHL{state}{S_3}$
			& $48$
			& $65$
			& $72$
			& $0.35$
			\\
			
			$\GTenToneHL{state}{S_4}$
			& $55$
			& $82$
			& $110$
			& $0.15$
			\\
			
			\midrule
			
			\GTenToneHL{expected}{EMV}
			& $\begin{array}{r}
				0.2(25)\\
				+0.3(35)\\
				+0.35(48)\\
				+0.15(55)\\
				\hline
				=40.55
			\end{array}$
			&
			$\begin{array}{r}
				0.2(5)\\
				+0.3(32)\\
				+0.35(65)\\
				+0.15(82)\\
				\hline
				=45.65
			\end{array}$
			&
			$\begin{array}{r}
				0.2(-25)\\
				+0.3(15)\\
				+0.35(72)\\
				+0.15(110)\\
				\hline
				=41.2
			\end{array}$
			&
			$\begin{array}{r}
				0.2\\
				+0.3\\
				+0.35\\
				+0.15\\
				\hline
				=1
			\end{array}$
			\\
			
			\bottomrule
		\end{tabular}
		
		\end{center}
	
	
	
\end{document}
